\question[4]
Konstruiere nach dem Verfahren aus dem Unterricht einen Huffman-Baum
für das Wort \texttt{kabelleger}. Gib für jeden Buchstaben die Codierung an.
Wieviele Bits werden gegenüber einer 8-Bit ASCII Codierung eingespart?

\begin{solutionbox}{12cm}
\begin{lstlisting}
Huffmanbaum für: kabelleger
...2 l
..3 #
...1 r
.6 #
..3 e
10 #
...1 g
..2 #
...1 b
.4 #
...1 a
..2 #
...1 k

Codierung:
k -> 000
a -> 001
b -> 010
e -> 10
l -> 111
g -> 011
r -> 110

kabelleger: 000 001 010 10 111 111 10 011 10 110
000001010101111111001110110  -  27 Bits
Ausgangstext mit 8 Bits pro Zeichen: 80 Bits
Einsparung: 53 Bits
\end{lstlisting}
\end{solutionbox}
